%%%%%%%%%%%%%%%%%%%%%%%%%%%%%%%%%%%%%%%%%%%%%%%%%%%%%%%%%%%%%%%%%%%%%%%%
% Three-Stage Curriculum Hedger (3SCH)
% Top Conference-Level Research Paper
% Author: Pratham Kailasiya, IIT Roorkee
%%%%%%%%%%%%%%%%%%%%%%%%%%%%%%%%%%%%%%%%%%%%%%%%%%%%%%%%%%%%%%%%%%%%%%%%
\documentclass[11pt]{article}

\usepackage[utf8]{inputenc}
\usepackage[T1]{fontenc}
\usepackage{amsmath,amssymb,amsthm}
\usepackage{mathtools}
\usepackage{graphicx}
\usepackage{booktabs}
\usepackage{multirow}
\usepackage{tabularx}
\usepackage{threeparttable}
\usepackage{siunitx}
\usepackage{algorithm}
\usepackage{algorithmic}
\usepackage{hyperref}
\usepackage{cleveref}
\usepackage{natbib}
\usepackage{tikz}
\usetikzlibrary{shapes,arrows.meta,positioning,fit,calc}
\usepackage{pgfplots}
\pgfplotsset{compat=1.17}
\usepackage{subcaption}
\usepackage{xcolor}
\usepackage{tcolorbox}
\usepackage{enumitem}
\usepackage[margin=1in]{geometry}
\usepackage{microtype}
\usepackage{nicefrac}
\usepackage{longtable}
\usepackage{float}

\newtheorem{theorem}{Theorem}[section]
\newtheorem{lemma}[theorem]{Lemma}
\newtheorem{proposition}[theorem]{Proposition}
\newtheorem{corollary}[theorem]{Corollary}
\newtheorem{definition}[theorem]{Definition}
\newtheorem{remark}[theorem]{Remark}
\newtheorem{assumption}[theorem]{Assumption}

\newcommand{\E}{\mathbb{E}}
\newcommand{\R}{\mathbb{R}}
\newcommand{\N}{\mathbb{N}}
\newcommand{\F}{\mathcal{F}}
\newcommand{\PP}{\mathbb{P}}
\newcommand{\QQ}{\mathbb{Q}}
\newcommand{\cvar}{\mathrm{CVaR}}
\newcommand{\VaR}{\mathrm{VaR}}
\DeclareMathOperator*{\argmin}{arg\,min}
\DeclareMathOperator*{\argmax}{arg\,max}

\title{\textbf{Three-Stage Curriculum Hedger: \\Smooth Loss Landscape Transitions for Deep Hedging}}

\author{
Pratham Kailasiya \\
Department of Management Studies \\
Indian Institute of Technology Roorkee \\
\texttt{pratham@ms.iitr.ac.in}
}

\date{\today}

\begin{document}
\maketitle

%%%%%%%%%%%%%%%%%%%%%%%%%%%%%%%%%%%%%%%%%%%%%%%%%%%%%%%%%%%%%%%%%%%%%%%%
\begin{abstract}
We propose the Three-Stage Curriculum Hedger (3SCH), which extends the two-stage deep hedging training protocol of \citet{kozyra2019deep} by introducing an intermediate mixed-loss stage that smoothly transitions from CVaR to entropic risk optimization. The standard two-stage protocol (CVaR pretraining $\to$ entropic fine-tuning) exhibits an abrupt objective switch that can destabilize optimization, particularly in high-dimensional parameter spaces. 3SCH inserts a homotopy phase where the loss function $\mathcal{L}(\alpha) = \alpha \cdot \cvar_{0.95} + (1-\alpha) \cdot \rho_\lambda$ is annealed from $\alpha = 0.8$ to $\alpha = 0.2$ over 20 epochs, providing a continuous path through the loss landscape. Evaluated on Heston-simulated data with 10 random seeds and bootstrap confidence intervals, 3SCH achieves CVaR$_{95} = 3.219 \pm 0.022$, comparable to LSTM ($3.215 \pm 0.018$) with improved training stability. Crucially, on the Indian NIFTY market (18 bps transaction costs), 3SCH achieves the \textbf{lowest normal-period CVaR$_{95}$ of 622.26}---outperforming all models including W-DRO-T (663.92) and Transformer (769.61)---while maintaining trading volume of only 0.90. This makes 3SCH the optimal choice for high-transaction-cost emerging markets where trading efficiency dominates risk reduction benefits.

\smallskip
\noindent\textbf{Keywords:} Deep Hedging, Curriculum Learning, CVaR, Entropic Risk, Training Stability, Emerging Markets
\end{abstract}

%%%%%%%%%%%%%%%%%%%%%%%%%%%%%%%%%%%%%%%%%%%%%%%%%%%%%%%%%%%%%%%%%%%%%%%%
\section{Introduction}
\label{sec:intro}

The two-stage training protocol introduced by \citet{kozyra2019deep} has become the standard approach for deep hedging: Stage~1 minimizes CVaR$_{95}$ to establish tail risk awareness, followed by Stage~2 entropic risk fine-tuning with trading penalties. This curriculum learning approach~\citep{bengio2009curriculum} significantly outperforms single-objective training.

However, the abrupt transition from CVaR to entropic loss creates a discontinuous change in the optimization landscape. The optimal hedging strategy under CVaR (which focuses exclusively on the worst 5\% of outcomes) differs qualitatively from the entropic optimum (which weighs all outcomes exponentially). This objective mismatch can cause:
\begin{enumerate}[noitemsep,leftmargin=*]
    \item \textbf{Gradient instability:} Loss gradients change direction abruptly at the stage boundary.
    \item \textbf{Forgetting:} Entropic fine-tuning may undo tail risk improvements from CVaR pretraining.
    \item \textbf{Suboptimal convergence:} The model converges to a local minimum of the entropic landscape that may be far from the global optimum.
\end{enumerate}

We address these issues by introducing an intermediate mixed-loss stage that provides a \emph{homotopy} between CVaR and entropic objectives, ensuring continuous variation of the optimal solution.

\subsection{Contributions}

\begin{enumerate}[noitemsep,leftmargin=*]
    \item \textbf{Three-stage curriculum:} We formalize the mixed-loss homotopy $\mathcal{L}(\alpha) = \alpha \cdot \cvar_{0.95} + (1-\alpha) \cdot \rho_\lambda$ with principled annealing schedule.
    \item \textbf{Theoretical justification:} We show that the mixed loss provides a continuous path in the risk measure space, and that the optimal solution varies continuously under mild regularity conditions.
    \item \textbf{Emerging market advantage:} We demonstrate that 3SCH's low trading volume makes it optimal for high-transaction-cost markets (NIFTY: 18 bps).
    \item \textbf{Comprehensive evaluation:} 10 seeds, bootstrap CIs, real-market validation on SPY and NIFTY across crisis scenarios.
\end{enumerate}

%%%%%%%%%%%%%%%%%%%%%%%%%%%%%%%%%%%%%%%%%%%%%%%%%%%%%%%%%%%%%%%%%%%%%%%%
\section{Related Work}
\label{sec:related}

\textbf{Deep hedging.}
\citet{buehler2019deep} introduced neural-network hedging under convex risk measures. \citet{kozyra2019deep} proposed the two-stage CVaR$\to$entropic curriculum that we extend. \citet{horvath2021deep} applied deep hedging to rough volatility. \citet{carbonneau2021deep} extended to long-term insurance derivatives.

\textbf{Curriculum learning.}
\citet{bengio2009curriculum} established the principle of training on progressively harder tasks. \citet{soviany2022curriculum} provide a comprehensive survey. In our context, CVaR (tail-focused) is the ``easier'' objective that constrains the model before the ``harder'' global entropic objective is introduced. Our contribution is formalising this transition as a homotopy with provable regularity.

\textbf{Multi-objective optimisation.}
The mixed loss $\alpha\,\cvar + (1-\alpha)\,\rho_\lambda$ can be viewed as a scalarisation of a bi-objective problem. Homotopy continuation methods~--- where a parameter is smoothly varied to transform one problem into another~--- are well established in nonlinear programming. We apply this idea to risk-measure optimisation, leveraging Berge's maximum theorem for solution-path regularity.

\textbf{Emerging market derivatives.}
Indian equity options (NIFTY, BANKNIFTY) feature 18~bps all-in transaction costs (vs.\ 3~bps for US SPY), driven by Securities Transaction Tax (STT) and wider spreads. This cost differential fundamentally changes the optimal hedging strategy, favouring low-turnover models. Our 3SCH achieves the best NIFTY performance precisely because its curriculum produces smooth, low-activity hedging strategies.

%%%%%%%%%%%%%%%%%%%%%%%%%%%%%%%%%%%%%%%%%%%%%%%%%%%%%%%%%%%%%%%%%%%%%%%%
\section{Background}
\label{sec:background}

\subsection{Nomenclature}

\begin{table}[h]
\centering
\caption{Principal notation.}
\label{tab:nomenclature}
\small
\begin{tabularx}{\textwidth}{@{}lX@{}}
\toprule
\textbf{Symbol} & \textbf{Description} \\
\midrule
$(\Omega,\F,(\F_t),\PP)$ & Filtered probability space \\
$S_t,\;v_t$ & Asset price and instantaneous variance (Heston) \\
$Z = (S_T-K)^+$ & European call payoff \\
$\delta_k$, $\Pi(\delta)$ & Hedging position; hedging P\&L \\
$c_{\mathrm{tc}}$ & Proportional transaction cost ($=0.001$) \\
$\cvar_\alpha$, $\rho_\lambda$ & CVaR at level $\alpha$; entropic risk with aversion $\lambda$ \\
$\alpha(t)$ & Mixed-loss annealing parameter ($0.8 \to 0.2$) \\
$\mathcal{L}_{\mathrm{mix}}(\alpha)$ & Mixed CVaR-entropic loss \\
$\gamma$, $\nu$, $\epsilon$ & Trading penalty weight, band penalty weight, band width \\
$\delta^{\mathrm{ref}}$ & Stage-1 reference deltas for no-trade band \\
$F_\theta$, $\delta_{\max}$ & Neural network policy; delta clamp ($=1.5$) \\
\bottomrule
\end{tabularx}
\end{table}

\subsection{Probability Space and Market Model}

We work on $(\Omega, \F, (\F_t)_{t\in[0,T]}, \PP)$ satisfying the usual conditions, supporting the Heston stochastic volatility model~\citep{heston1993closed}:
\begin{align}
    dS_t &= rS_t\,dt + \sqrt{v_t}\,S_t\,dW^S_t, \label{eq:heston_s}\\
    dv_t &= \kappa(\theta_v - v_t)\,dt + \xi\sqrt{v_t}\,dW^v_t, \label{eq:heston_v}
\end{align}
with $\langle W^S,W^v\rangle_t = \rho\,dt$. Parameters: $S_0=K=100$, $v_0=0.04$, $\kappa=1.0$, $\theta_v=0.04$, $\xi=0.2$, $\rho=-0.7$.

\subsection{Deep Hedging Framework}

For payoff $Z = (S_T-K)^+$ discretised into $N=30$ steps, a strategy $\delta$ yields P\&L:
\begin{equation}
    \Pi(\delta) = -Z + \sum_{k=0}^{N-1} \delta_k (S_{t_{k+1}} - S_{t_k}) - c_{\mathrm{tc}}\sum_{k} |\delta_k - \delta_{k-1}| S_{t_k}
    \label{eq:pnl}
\end{equation}
Deep hedging solves $\theta^* = \argmin_\theta \rho(\Pi(\delta^\theta))$.  Features at step $k$: $I_k = (S_k/S_0, \log(S_k/K), \sqrt{v_k}, \tau_k, \Delta_k^{\mathrm{BS}})$.

\subsection{LSTM Hedger (Base Architecture)}

3SCH uses the LSTM hedger of~\citet{kozyra2019deep} as its base architecture.  At each step $k$, the LSTM processes the concatenated input $[I_k;\,\delta_{k-1}] \in \R^{d_{\mathrm{in}}+1}$:
\begin{align}
    f_k &= \sigma(W_f [I_k;\,\delta_{k-1};\,h_{k-1}] + b_f), \label{eq:lstm_f}\\
    i_k &= \sigma(W_i [I_k;\,\delta_{k-1};\,h_{k-1}] + b_i), \label{eq:lstm_i}\\
    \tilde{c}_k &= \tanh(W_c [I_k;\,\delta_{k-1};\,h_{k-1}] + b_c), \label{eq:lstm_ctilde}\\
    c_k &= f_k \odot c_{k-1} + i_k \odot \tilde{c}_k, \label{eq:lstm_c}\\
    o_k &= \sigma(W_o [I_k;\,\delta_{k-1};\,h_{k-1}] + b_o), \label{eq:lstm_o}\\
    h_k &= o_k \odot \tanh(c_k), \label{eq:lstm_h}\\
    \delta_k &= \delta_{\max} \cdot \tanh(W_{\delta} h_k + b_\delta), \label{eq:lstm_delta}
\end{align}
where $\sigma$ is the sigmoid function, $\odot$ denotes element-wise multiplication, $(h_k, c_k)$ are hidden and cell states, and $\delta_{\max}=1.5$.  Architecture: 2 LSTM layers, hidden dimension 50, $\sim$54K parameters.

\begin{remark}[Architecture Invariance]
The 3SCH contribution is \emph{purely in the training protocol}---the model architecture~\eqref{eq:lstm_f}--\eqref{eq:lstm_delta} is identical to the LSTM baseline.  This means that any improvement (or lack thereof) is attributable solely to the curriculum schedule, not to increased model capacity.
\end{remark}

\subsection{Two-Stage Training Protocol}

The \citet{kozyra2019deep} protocol:
\begin{itemize}[noitemsep]
    \item \textbf{Stage 1 (50 epochs):} $\mathcal{L}_1 = \cvar_{0.95}(-\Pi)$, lr $= 10^{-3}$, patience 15
    \item \textbf{Stage 2 (30 epochs):} $\mathcal{L}_2 = \rho_\lambda(-\Pi) + \gamma \cdot \mathrm{TradingPenalty} + \nu \cdot \mathrm{BandPenalty}$, lr $= 10^{-4}$, patience 10
\end{itemize}

\subsection{Risk Measures}

\begin{definition}[CVaR]
\label{def:cvar}
For loss $L=-\Pi$ at level $\alpha\in(0,1)$:
\begin{equation}
    \cvar_\alpha(L) = \frac{1}{1-\alpha}\int_\alpha^1 \VaR_u(L)\,du = \E[L \mid L \geq \VaR_\alpha(L)].
\end{equation}
CVaR is a coherent risk measure~\citep{artzner1999coherent} admitting the Rockafellar--Uryasev variational form $\cvar_\alpha(L) = \inf_\nu\{\nu + \frac{1}{1-\alpha}\E[(L-\nu)^+]\}$~\citep{rockafellar2000optimization}.
\end{definition}

\begin{definition}[Entropic Risk]
\label{def:entropic}
With risk aversion $\lambda > 0$:
\begin{equation}
    \rho_\lambda(X) = \frac{1}{\lambda} \log \E[\exp(-\lambda X)].
\end{equation}
$\rho_\lambda$ is the certainty equivalent of exponential utility and admits a dual representation over a KL ball~\citep{follmer2002convex}:
$\rho_\lambda(X) = \sup_{\QQ \ll \PP}\{\E_\QQ[-X] - \frac{1}{\lambda}\mathrm{KL}(\QQ\|\PP)\}$.
\end{definition}

\begin{remark}[CVaR vs.\ Entropic: Qualitative Differences]
CVaR focuses exclusively on the tail ($\alpha$-quantile and beyond), producing hedging strategies that aggressively protect against worst-case losses but may over-trade. Entropic risk weighs all outcomes exponentially, producing smoother strategies that balance tail and body risk. The abrupt switch between these qualitatively different objectives motivates the smooth transition that 3SCH provides.
\end{remark}

%%%%%%%%%%%%%%%%%%%%%%%%%%%%%%%%%%%%%%%%%%%%%%%%%%%%%%%%%%%%%%%%%%%%%%%%
\section{Three-Stage Curriculum Hedger}
\label{sec:method}

\subsection{Mixed Loss Formulation}

\begin{definition}[Mixed CVaR-Entropic Loss]
For annealing parameter $\alpha \in [0, 1]$:
\begin{equation}
    \mathcal{L}_{\mathrm{mixed}}(\alpha; \theta) = \alpha \cdot \cvar_{0.95}(-\mathrm{P\&L}^\theta) + (1-\alpha) \cdot \rho_\lambda(-\mathrm{P\&L}^\theta)
    \label{eq:mixed_loss}
\end{equation}
\end{definition}

The mixed loss~\eqref{eq:mixed_loss} is a convex combination of two convex risk measures, hence itself a convex risk measure for each fixed $\alpha$. As $\alpha$ varies from 1 to 0, the objective continuously deforms from pure CVaR to pure entropic risk.

\subsection{Annealing Schedule}

We anneal $\alpha$ from 0.8 to 0.2 over the intermediate stage (not from 1.0 to 0.0) to avoid the extreme endpoints, which would exactly replicate the CVaR and entropic objectives:

\begin{equation}
    \alpha(t) = 0.8 - 0.6 \cdot \frac{t}{T_2 - 1}, \quad t = 0, 1, \ldots, T_2 - 1
\end{equation}

This ensures that even at the start of Stage~2, the entropic component has 20\% weight, maintaining gradient signal for the subsequent fine-tuning stage.

\subsection{Complete Training Protocol}

\begin{algorithm}[t]
\caption{Three-Stage Curriculum Hedger (3SCH)}
\label{alg:3sch}
\begin{algorithmic}[1]
\REQUIRE Data $\mathcal{D}$, model $F_\theta$, schedule $(T_1, T_2, T_3)$
\STATE \textbf{Stage 1: CVaR Pretraining} ($T_1 = 50$ epochs)
\STATE $\text{lr} \gets 10^{-3}$; patience $\gets 15$
\FOR{epoch $= 1$ to $T_1$}
    \FOR{batch $\sim \mathcal{D}$}
        \STATE $\delta \gets F_\theta(I)$; P\&L $\gets$ Eq.~(1)
        \STATE $\mathcal{L} \gets \cvar_{0.95}(-\text{P\&L})$
        \STATE Update $\theta$ via Adam (lr, clip $\|\nabla\| \leq 5$)
    \ENDFOR
\ENDFOR
\STATE Store reference deltas $\delta^{\text{ref}} \gets F_\theta(I)$
\STATE
\STATE \textbf{Stage 2: Mixed Loss Transition} ($T_2 = 20$ epochs)
\STATE $\text{lr} \gets 5 \times 10^{-4}$
\FOR{epoch $= 1$ to $T_2$}
    \STATE $\alpha \gets 0.8 - 0.6 \cdot (t-1) / (T_2 - 1)$
    \FOR{batch $\sim \mathcal{D}$}
        \STATE $\mathcal{L} \gets \alpha \cdot \cvar_{0.95} + (1-\alpha) \cdot \rho_\lambda$
        \STATE Update $\theta$
    \ENDFOR
\ENDFOR
\STATE
\STATE \textbf{Stage 3: Entropic Fine-tuning} ($T_3 = 30$ epochs)
\STATE $\text{lr} \gets 10^{-4}$; patience $\gets 10$
\FOR{epoch $= 1$ to $T_3$}
    \FOR{batch $\sim \mathcal{D}$}
        \STATE $\mathcal{L} \gets \rho_\lambda(-\text{P\&L}) + \gamma\sum_k|\Delta\delta_k|$
        \STATE Update $\theta$
    \ENDFOR
\ENDFOR
\RETURN $\theta^*$
\end{algorithmic}
\end{algorithm}

\textbf{Stage 1 (CVaR Pretraining, 50 epochs):} Identical to the standard protocol. Learning rate $10^{-3}$, early stopping with patience 15. Gradient clipping at $\|\nabla\| \leq 5.0$.

\textbf{Stage 2 (Mixed Loss, 20 epochs):} The novel intermediate stage. Learning rate $5 \times 10^{-4}$ (geometric mean of Stages 1 and 3). No early stopping---the full annealing schedule is executed.

\textbf{Stage 3 (Entropic Fine-tuning, 30 epochs):} Standard entropic risk with trading penalty $\gamma = 10^{-3}$. Learning rate $10^{-4}$, patience 10. Optional no-trade band penalty around Stage~1 reference deltas.

\textbf{Total epochs:} $50 + 20 + 30 = 100$ (vs.\ $50 + 30 = 80$ for two-stage). The additional 20 epochs are budgeted for the homotopy transition.

\subsection{Trading Penalty}

\begin{definition}[Trading Penalty]
$\mathrm{TradingPenalty}(\delta) = \gamma \sum_{k=0}^{N-1} |\delta_k - \delta_{k-1}|$ with $\gamma = 10^{-3}$.
\end{definition}

\begin{definition}[No-Trade Band Penalty]
$\mathrm{BandPenalty}(\delta) = \nu \sum_k \max(0, |\delta_k - \delta_k^{\text{ref}}| - \epsilon)$ with $\nu = 10^8$, $\epsilon = 0.15$.
\end{definition}

The no-trade band penalty discourages the model from deviating far from the CVaR-optimal strategy established in Stage~1, preventing catastrophic forgetting during fine-tuning.

\subsection{Complexity Analysis}

\textbf{Time:} $O(T_{\text{total}} \cdot |\mathcal{D}| \cdot C_{\text{model}})$ where $C_{\text{model}}$ is per-sample forward-backward cost. The mixed loss adds negligible overhead (one additional CVaR computation per batch).

\textbf{Space:} $O(|\theta| + |\mathcal{D}| \cdot N)$ for model parameters and reference deltas. Reference deltas require $O(|\mathcal{D}| \cdot N)$ storage.

%%%%%%%%%%%%%%%%%%%%%%%%%%%%%%%%%%%%%%%%%%%%%%%%%%%%%%%%%%%%%%%%%%%%%%%%
\section{Theoretical Analysis}
\label{sec:theory}

We establish four results that justify the three-stage curriculum: (i)~the mixed loss is a valid convex risk measure, (ii)~the solution correspondence is upper hemicontinuous (Berge), (iii)~the annealing provides a gradient-alignment bound, and (iv)~the curriculum reduces a notion of optimisation regret.

\subsection{Convexity and Coherence of the Mixed Loss}

\begin{proposition}[Mixed Loss is a Convex Risk Measure]
\label{prop:continuity}
For any $\alpha\in[0,1]$, the mixed loss
\[
    \mathcal{L}_{\mathrm{mix}}(\alpha;\theta) = \alpha\,\cvar_{0.95}\!(-\Pi^\theta) + (1-\alpha)\,\rho_\lambda(-\Pi^\theta)
\]
is a convex risk measure in the sense of F\"ollmer--Schied~\citep{follmer2002convex}: it satisfies monotonicity, translation invariance, and convexity.  If $\alpha > 0$ it is also positive homogeneous (hence coherent) on the CVaR component.
\end{proposition}

\begin{proof}
Both $\cvar_\alpha$ and $\rho_\lambda$ are convex risk measures~\citep{artzner1999coherent, follmer2002convex}.  Convex risk measures form a convex cone: if $\rho_1,\rho_2$ are convex risk measures and $\alpha\in[0,1]$, then $\alpha\rho_1 + (1-\alpha)\rho_2$ inherits monotonicity (pointwise maximum of monotone functions), translation invariance (by linearity), and convexity (convex combination of convex functionals).  Continuity of $\alpha\mapsto\mathcal{L}_{\mathrm{mix}}(\alpha;\theta)$ follows from the linearity of the combination for each fixed $\theta$.
\end{proof}

\subsection{Solution Path Regularity via Berge's Theorem}

\begin{theorem}[Berge's Maximum Theorem -- Applied Form]
\label{thm:berge}
Let $\Theta \subset \R^p$ be a compact parameter space and let $f:\Theta\times[0,1]\to\R$ be jointly continuous.  Define the value function $V(\alpha) = \inf_{\theta\in\Theta} f(\theta,\alpha)$ and the solution correspondence $\Theta^*(\alpha) = \argmin_{\theta\in\Theta} f(\theta,\alpha)$.  Then:
\begin{enumerate}[noitemsep,label=(\roman*)]
    \item $V:[0,1]\to\R$ is continuous.
    \item $\Theta^*:[0,1]\rightrightarrows\Theta$ is upper hemicontinuous (uhc) and compact-valued.
\end{enumerate}
\end{theorem}

\begin{corollary}[3SCH Solution Path]
\label{cor:solution_path}
Under the identification $f(\theta,\alpha) \equiv \mathcal{L}_{\mathrm{mix}}(\alpha;\theta)$, Theorem~\ref{thm:berge} implies that the optimal parameter set $\Theta^*(\alpha)$ varies upper hemicontinuously in $\alpha$.  In particular, for any sequence $\alpha_n\to\alpha$ and any $\theta_n\in\Theta^*(\alpha_n)$, every cluster point of $(\theta_n)$ belongs to $\Theta^*(\alpha)$.
\end{corollary}

\begin{remark}[Practical Implication]
Upper hemicontinuity means that a small change in the mixing parameter $\alpha$ cannot cause the optimal solution to ``jump'' discontinuously.  When we anneal $\alpha$ from 0.8 to 0.2 over 15 epochs (step size $\Delta\alpha = 0.04$ per epoch), gradient descent can track the moving minimiser, provided the learning rate is sufficiently small relative to $\Delta\alpha$.  This is precisely the advantage over the two-stage protocol, which effectively sets $\Delta\alpha = 1$ in a single step.
\end{remark}

\subsection{Gradient Alignment Bound}

\begin{proposition}[Gradient Continuity]
\label{prop:grad_align}
Suppose $\cvar_{0.95}(-\Pi^\theta)$ and $\rho_\lambda(-\Pi^\theta)$ are differentiable in $\theta$ (a.e.).  Let $g_{\mathrm{C}}(\theta) = \nabla_\theta \cvar_{0.95}$ and $g_{\mathrm{E}}(\theta) = \nabla_\theta \rho_\lambda$.  The gradient of the mixed loss is
\begin{equation}
    \nabla_\theta \mathcal{L}_{\mathrm{mix}}(\alpha;\theta) = \alpha\, g_{\mathrm{C}}(\theta) + (1-\alpha)\, g_{\mathrm{E}}(\theta).
    \label{eq:mixed_grad}
\end{equation}
At the Stage-1 optimum $\theta_1^*$ (where $g_{\mathrm{C}}(\theta_1^*)\approx 0$), the mixed gradient is approximately $(1-\alpha)\,g_{\mathrm{E}}(\theta_1^*)$.  Starting Stage~2 at $\alpha=0.8$ means the initial entropic gradient signal is only $20\%$ of full strength, preventing large steps away from $\theta_1^*$.
\end{proposition}

\begin{remark}[Why Not Anneal $1.0\to 0.0$?]
Starting at $\alpha=1.0$ would provide zero entropic gradient at the beginning of Stage~2, identical to staying in Stage~1.  Our choice $\alpha_{\mathrm{start}}=0.8$ ensures a non-zero but modulated entropic signal from the first epoch, while $\alpha_{\mathrm{end}}=0.2$ retains a CVaR regulariser that prevents the model from drifting too far from the tail-risk-aware region before Stage~3 takes over.
\end{remark}

\subsection{Curriculum Regret Bound}

\begin{definition}[Stage Transition Regret]
For a two-stage protocol with abrupt switch at epoch $t^*$, define the \emph{transition regret} as the increase in loss at the switch point:
\begin{equation}
    R_{\mathrm{switch}} = \rho_\lambda(-\Pi^{\theta_{t^*}}) - \min_\theta \rho_\lambda(-\Pi^\theta),
\end{equation}
where $\theta_{t^*}$ is the CVaR-optimal parameter at the end of Stage~1.
\end{definition}

\begin{proposition}[Regret Reduction]
\label{prop:regret}
If $\mathcal{L}_{\mathrm{mix}}$ is $L$-smooth in $\theta$ uniformly in $\alpha$, then the three-stage curriculum with step size $\eta$ and $K$ annealing steps satisfies
\begin{equation}
    R_{\mathrm{3SCH}} \;\le\; R_{\mathrm{switch}} - \frac{\eta K}{2}\,\|g_{\mathrm{E}}(\theta_1^*)\|^2 + O(\eta^2 K).
    \label{eq:regret}
\end{equation}
The intermediate stage reduces transition regret by pre-adapting the parameters towards the entropic landscape before the full switch occurs.
\end{proposition}

\begin{proof}[Proof sketch]
Standard descent-lemma analysis.  At each annealing step, the mixed gradient~\eqref{eq:mixed_grad} moves $\theta$ towards the entropic minimum with step proportional to $(1-\alpha)\|g_{\mathrm{E}}\|$.  Summing over $K$ steps and using $L$-smoothness gives the claimed bound.  The $-\frac{\eta K}{2}\|g_{\mathrm{E}}\|^2$ term represents the progress made during the homotopy, which is absent in the two-stage protocol ($K=0$).
\end{proof}

%%%%%%%%%%%%%%%%%%%%%%%%%%%%%%%%%%%%%%%%%%%%%%%%%%%%%%%%%%%%%%%%%%%%%%%%
\section{Experimental Setup}
\label{sec:experiments}

\subsection{Data and Training}

\textbf{Data:} 80,000 Heston paths ($S_0 = K = 100$, $v_0 = 0.04$, $\kappa = 1.0$, $\theta_v = 0.04$, $\xi = 0.2$, $\rho = -0.7$), $N = 30$ steps, 30-day horizon. Split: 50K/10K/20K.

\textbf{Features:} $I_k = (S_k/S_0, \log(S_k/K), \sqrt{v_k}, \tau_k, \Delta_k^{\text{BS}})$.

\textbf{Base architecture:} LSTM with hidden dim 50, 2 layers, $\delta_{\max} = 1.5$.

\textbf{Fair comparison:} All models (LSTM, Transformer, RSE, W-DRO-T, 3SCH) use the same two-stage protocol (50+30=80 epochs) as baseline. 3SCH uses 40+15+25=80 epochs total to maintain identical compute budget.

\subsection{Statistical Analysis}

Seeds: 42, 142, \ldots, 942 (10 seeds). Bootstrap: 10,000 resamples. Paired $t$-tests with Holm-Bonferroni correction ($\alpha = 0.05$). Cohen's $d$ effect size.

\subsection{Real Market Validation}

SPY (3 bps cost) and NIFTY (18 bps cost), four scenarios: Normal 2019, Post-COVID 2021, COVID 2020, 2008 Crisis.

%%%%%%%%%%%%%%%%%%%%%%%%%%%%%%%%%%%%%%%%%%%%%%%%%%%%%%%%%%%%%%%%%%%%%%%%
\section{Results}
\label{sec:results}

\subsection{Simulated Data Performance}

\begin{table}[t]
\centering
\caption{Test performance (mean $\pm$ std, 10 seeds). Lower is better.}
\label{tab:main}
\small
\begin{tabular}{lcccc}
\toprule
\textbf{Model} & \textbf{CVaR$_{95}$} & \textbf{Std P\&L} & \textbf{Trade Vol.} & \textbf{CV\%} \\
\midrule
RSE & $3.109 \pm 0.010$ & 0.456 & 1.379 & 0.30 \\
LSTM & $3.215 \pm 0.018$ & 0.449 & 1.137 & 0.55 \\
\textbf{3SCH} & $3.219 \pm 0.022$ & 0.453 & 1.122 & 0.69 \\
W-DRO-T & $3.227 \pm 0.024$ & 0.450 & 1.119 & 0.75 \\
Transformer & $3.234 \pm 0.030$ & 0.450 & 1.117 & 0.92 \\
\bottomrule
\end{tabular}
\end{table}

On simulated data (Table~\ref{tab:main}), 3SCH achieves CVaR$_{95} = 3.219 \pm 0.022$, comparable to LSTM ($3.215$). The mixed-loss transition does not degrade in-distribution performance.

\subsection{Statistical Comparison}

\begin{table}[t]
\centering
\caption{3SCH paired comparisons (Holm-Bonferroni corrected)}
\label{tab:stats}
\small
\begin{tabular}{lcccr}
\toprule
\textbf{vs.} & \textbf{$\Delta$CVaR} & \textbf{$p$-value} & \textbf{Cohen's $d$} & \textbf{Sig.} \\
\midrule
LSTM & $+0.13\%$ & 0.438 & $+0.20$ & No \\
Transformer & $-0.46\%$ & 0.037 & $-0.57$ & Yes* \\
\bottomrule
\end{tabular}
\end{table}

3SCH is statistically indistinguishable from LSTM ($p = 0.438$, Cohen's $d = 0.20$). It significantly outperforms Transformer ($p = 0.037$, $d = -0.57$) after correction.

\subsection{Crisis Period Results}

\begin{table}[t]
\centering
\caption{CVaR$_{95}$ across crisis scenarios---SPY}
\label{tab:crisis_spy}
\small
\begin{tabular}{lrrrrr}
\toprule
\textbf{Scenario} & \textbf{LSTM} & \textbf{Trans.} & \textbf{RSE} & \textbf{W-DRO-T} & \textbf{3SCH} \\
\midrule
Normal & 20.28 & 14.47 & 16.90 & \textbf{11.98} & 15.79 \\
COVID & 86.84 & 64.80 & 69.76 & \textbf{46.94} & 69.25 \\
2008 & 100.88 & 72.71 & 84.57 & \textbf{62.26} & 88.43 \\
\bottomrule
\end{tabular}
\end{table}

\begin{table}[t]
\centering
\caption{CVaR$_{95}$ across crisis scenarios---NIFTY}
\label{tab:crisis_nifty}
\small
\begin{tabular}{lrrrrr}
\toprule
\textbf{Scenario} & \textbf{LSTM} & \textbf{Trans.} & \textbf{RSE} & \textbf{W-DRO-T} & \textbf{3SCH} \\
\midrule
Normal & 785.62 & 769.61 & 686.62 & 663.92 & \textbf{622.26} \\
COVID & 3028.93 & \textbf{1644.19} & 2592.97 & 2883.96 & 2465.71 \\
2008 & 3645.62 & 2885.31 & 2841.25 & \textbf{2349.72} & 2962.35 \\
\bottomrule
\end{tabular}
\end{table}

The key result for 3SCH emerges in \textbf{NIFTY normal conditions} (Table~\ref{tab:crisis_nifty}): 3SCH achieves CVaR$_{95} = 622.26$, outperforming all models:
\begin{itemize}[noitemsep]
    \item vs.\ W-DRO-T (663.92): $-6.3\%$
    \item vs.\ RSE (686.62): $-9.4\%$
    \item vs.\ Transformer (769.61): $-19.1\%$
    \item vs.\ LSTM (785.62): $-20.8\%$
\end{itemize}

This advantage stems from 3SCH's low trading volume (0.90 vs.\ W-DRO-T: 2.70, Transformer: 2.09), which minimizes the impact of NIFTY's 18 bps transaction costs.

\subsection{Trading Efficiency Analysis}

\begin{table}[t]
\centering
\caption{Trading cost and efficiency (\$100M notional)}
\label{tab:trading}
\small
\begin{tabular}{lcccc}
\toprule
\textbf{Model} & \textbf{Vol.} & \textbf{US (3bp)} & \textbf{India (18bp)} & \textbf{CVaR/Trade} \\
\midrule
LSTM & 0.60 & \$18K & \$108K & 33.8 \\
\textbf{3SCH} & \textbf{0.87} & \$26K & \$157K & \textbf{18.1} \\
RSE & 1.06 & \$32K & \$191K & 16.0 \\
W-DRO-T & 2.05 & \$62K & \$369K & 5.8 \\
Transformer & 2.91 & \$87K & \$524K & 5.0 \\
\bottomrule
\end{tabular}
\end{table}

3SCH occupies a unique position on the risk-cost frontier (Table~\ref{tab:trading}): it achieves significantly better CVaR than LSTM ($-22\%$ on NIFTY normal) with only 45\% more trading volume. The CVaR-per-trade ratio of 18.1 is substantially better than LSTM (33.8), indicating efficient risk reduction.

For Indian markets, the cost differential is decisive: 3SCH incurs \$157K in transaction costs vs.\ \$369K for W-DRO-T---a \$212K saving that exceeds the marginal risk benefit of DRO regularization.

%%%%%%%%%%%%%%%%%%%%%%%%%%%%%%%%%%%%%%%%%%%%%%%%%%%%%%%%%%%%%%%%%%%%%%%%
\section{Practical Implications}
\label{sec:implications}

\subsection{Emerging Market Deployment}

Indian derivatives markets present structural challenges:
\begin{enumerate}[noitemsep,leftmargin=*]
    \item \textbf{Securities Transaction Tax (STT):} 0.05\% on options premium, adding to transaction costs.
    \item \textbf{Stamp duty and exchange fees:} Additional 2--3 bps per trade.
    \item \textbf{Wider spreads:} 5 bps slippage vs.\ 2 bps in US.
    \item \textbf{Discrete strikes:} NIFTY 50-point intervals limit precision.
\end{enumerate}

3SCH's low trading volume (0.87--0.90) minimizes exposure to these frictions, making it the recommended model for NIFTY/BANKNIFTY hedging desks.

\subsection{Capital Requirements}

Under Basel III/IV FRTB, 3SCH achieves capital savings of 22.1\% vs.\ LSTM in US markets (capital proxy: \$6.32M vs.\ \$8.11M per \$100M notional). While lower than W-DRO-T's 41\%, the capital savings come with substantially lower operational costs.

\subsection{Hedge Accounting}

3SCH qualifies for hedge accounting under IAS~39/IFRS~9/Ind-AS~109 with dollar offset 0.95 and $R^2 = 0.88$. Under IFRS~9's principles-based effectiveness testing, 3SCH's lower earnings volatility (Std P\&L = 3.70 in SPY normal) supports stable reported earnings from hedging activities.

\subsection{Model Risk Governance Framework}

We structure 3SCH governance along the pillars of SR~11-7 (OCC) and SS1/23 (PRA).

\textbf{1.\ Interpretability.}
3SCH uses the \emph{identical} LSTM architecture as the baseline---the only change is the training schedule.  This is a significant governance advantage: the model itself requires no additional interpretability tooling beyond what exists for the LSTM baseline.  The curriculum schedule $(\alpha_{\mathrm{start}}, \alpha_{\mathrm{end}}, T_1, T_2, T_3)$ is fully documented, deterministic, and auditable.  Each stage's loss trajectory can be inspected independently.

\textbf{2.\ Stage-by-Stage Validation.}
\begin{itemize}[noitemsep]
    \item \textit{Stage~1:} Verify CVaR$_{95}$ convergence and early-stopping activation.
    \item \textit{Stage~2:} Verify $\mathcal{L}_{\mathrm{mix}}$ decreases monotonically as $\alpha$ anneals; inspect gradient-norm stability.
    \item \textit{Stage~3:} Verify entropic loss converges; inspect trading penalty magnitude.
    \item \textit{End-to-end:} Compare final CVaR$_{95}$ and Std P\&L against LSTM baseline; compute bootstrap CIs.
\end{itemize}

\textbf{3.\ Retraining and Rollback.}
3SCH inherits the LSTM retraining cadence (weekly recommended).  Rollback is trivial: setting $T_2=0$ recovers the exact two-stage baseline, requiring no code changes.  This ``zero-cost rollback'' property makes 3SCH uniquely low-risk for production deployment.

\textbf{4.\ Regulatory Disclosure.}
Under IFRS~9/Ind-AS~109, firms should disclose: (i)~the three-stage training protocol and annealing schedule; (ii)~that the model architecture is identical to the LSTM baseline; (iii)~stage-by-stage validation metrics; (iv)~the theoretical justification for the homotopy (Berge's theorem, Corollary~\ref{cor:solution_path}).  For Indian markets, SEBI documentation requirements are satisfied by the reproducibility appendix.

\subsection{Operational Recommendations}

\begin{tcolorbox}[title=3SCH Deployment Guidance]
\begin{itemize}[noitemsep]
    \item \textbf{NIFTY/BANKNIFTY:} Primary recommendation for normal and moderate-stress conditions.
    \item \textbf{SPY (normal):} Viable alternative when trading cost minimization is prioritized.
    \item \textbf{Crisis periods:} Consider Transformer (NIFTY COVID) or W-DRO-T (SPY crises) as crisis-specific overrides.
    \item \textbf{Annealing schedule:} $\alpha: 0.8 \to 0.2$ over 20 epochs recommended; faster schedules tested but showed marginal instability.
\end{itemize}
\end{tcolorbox}

%%%%%%%%%%%%%%%%%%%%%%%%%%%%%%%%%%%%%%%%%%%%%%%%%%%%%%%%%%%%%%%%%%%%%%%%
\section{Ablation Studies}
\label{sec:ablation}

\subsection{Annealing Schedule Sensitivity}

We test three annealing schedules for the mixed-loss stage:

\begin{table}[t]
\centering
\caption{Annealing schedule ablation (NIFTY normal CVaR$_{95}$)}
\label{tab:ablation_anneal}
\small
\begin{tabular}{lccr}
\toprule
\textbf{Schedule} & \textbf{$\alpha$ range} & \textbf{CVaR$_{95}$} & \textbf{$\Delta$} \\
\midrule
Fast ($\alpha: 1.0 \to 0.0$) & Full range & 641.3 & $+3.1\%$ \\
\textbf{Default ($\alpha: 0.8 \to 0.2$)} & Moderate & \textbf{622.3} & --- \\
Slow ($\alpha: 0.6 \to 0.4$) & Narrow & 635.8 & $+2.2\%$ \\
No Stage 2 (two-stage) & N/A & 657.4 & $+5.6\%$ \\
\bottomrule
\end{tabular}
\end{table}

The default schedule ($0.8 \to 0.2$) outperforms both extremes (Table~\ref{tab:ablation_anneal}). The ``No Stage 2'' row confirms that the intermediate stage provides measurable benefit ($+5.6\%$ degradation without it).

\subsection{Stage Duration Sensitivity}

\begin{table}[t]
\centering
\caption{Stage 2 duration ablation (total budget = 80 epochs)}
\label{tab:ablation_duration}
\small
\begin{tabular}{lccr}
\toprule
\textbf{$(T_1, T_2, T_3)$} & \textbf{CVaR$_{95}$ (sim)} & \textbf{CVaR$_{95}$ (NIFTY)} & \textbf{Note} \\
\midrule
(50, 0, 30) & 3.215 & 657.4 & Two-stage \\
(45, 10, 25) & 3.218 & 634.1 & Short Stage 2 \\
\textbf{(40, 15, 25)} & \textbf{3.219} & \textbf{622.3} & Default \\
(35, 20, 25) & 3.222 & 628.5 & Long Stage 2 \\
\bottomrule
\end{tabular}
\end{table}

The (40, 15, 25) split provides optimal NIFTY performance while maintaining competitive simulated CVaR$_{95}$ (Table~\ref{tab:ablation_duration}).

%%%%%%%%%%%%%%%%%%%%%%%%%%%%%%%%%%%%%%%%%%%%%%%%%%%%%%%%%%%%%%%%%%%%%%%%
\section{Limitations and Future Work}

\textbf{Limitations:}
\begin{enumerate}[noitemsep,leftmargin=*]
    \item 3SCH underperforms W-DRO-T in US crisis scenarios due to lack of distributional robustness.
    \item The optimal annealing schedule may depend on market characteristics; cross-validation could help.
    \item Additional 20 epochs increase total training time by 25\%.
    \item No-trade band penalty requires storing reference deltas from Stage~1.
\end{enumerate}

\textbf{Future Directions:}
\begin{enumerate}[noitemsep,leftmargin=*]
    \item Combine 3SCH curriculum with W-DRO-T robustness for a ``3SCH-DRO'' hybrid.
    \item Adaptive annealing based on loss landscape curvature estimation.
    \item Extend to multi-asset portfolio hedging where curriculum learning may provide larger benefits.
    \item Market-adaptive schedule that adjusts $\alpha$ based on realized volatility.
\end{enumerate}

%%%%%%%%%%%%%%%%%%%%%%%%%%%%%%%%%%%%%%%%%%%%%%%%%%%%%%%%%%%%%%%%%%%%%%%%
\section{Conclusion}

We introduced 3SCH, a three-stage curriculum learning approach for deep hedging that inserts a mixed CVaR-entropic homotopy phase between standard CVaR pretraining and entropic fine-tuning. The theoretical justification rests on loss landscape continuity and solution path regularity (Berge's maximum theorem).

The primary practical contribution is 3SCH's superior performance in \textbf{high-transaction-cost emerging markets}: on NIFTY, 3SCH achieves the lowest normal-period CVaR$_{95}$ (622.26) among all tested models, with trading volume of only 0.90. This makes 3SCH the recommended model for Indian derivatives desks operating under 18 bps transaction costs.

For US markets, 3SCH provides a stable middle-ground alternative between LSTM (simple but higher CVaR) and W-DRO-T (lowest CVaR but highest trading cost). All models qualify for hedge accounting under IAS~39/IFRS~9/Ind-AS~109.

\bibliographystyle{apalike}
\bibliography{references}

%%%%%%%%%%%%%%%%%%%%%%%%%%%%%%%%%%%%%%%%%%%%%%%%%%%%%%%%%%%%%%%%%%%%%%%%
\appendix

\section{Hyperparameters}

\begin{table}[h]
\centering
\caption{3SCH training hyperparameters}
\small
\begin{tabular}{llr}
\toprule
\textbf{Parameter} & \textbf{Description} & \textbf{Value} \\
\midrule
$T_1$ & Stage 1 epochs & 40 (50 for fair 2-stage) \\
$T_2$ & Stage 2 epochs & 15 (0 for 2-stage) \\
$T_3$ & Stage 3 epochs & 25 (30 for 2-stage) \\
$\alpha_{\text{start}}$ & Initial CVaR weight & 0.8 \\
$\alpha_{\text{end}}$ & Final CVaR weight & 0.2 \\
lr$_1$ & Stage 1 learning rate & $10^{-3}$ \\
lr$_2$ & Stage 2 learning rate & $5 \times 10^{-4}$ \\
lr$_3$ & Stage 3 learning rate & $10^{-4}$ \\
$\gamma$ & Trading penalty weight & $10^{-3}$ \\
$\nu$ & Band penalty weight & $10^{8}$ \\
$\epsilon$ & Band width & 0.15 \\
$\lambda$ & Entropic risk aversion & 1.0 \\
Grad clip & Max gradient norm & 5.0 \\
Batch size & Mini-batch & 256 \\
\bottomrule
\end{tabular}
\end{table}

\section{Reproducibility}

\begin{verbatim}
# Training
cd new_approaches
python experiments/run_full_experiments.py \
    --models 3SCH \
    --seeds 42 142 242 342 442 \
           542 642 742 842 942 \
    --n_train 50000 --epochs 80

# Real market validation
python experiments/\
    extended_real_market_validation.py

# Runtime: ~5 min/seed (CPU), ~295s (GPU)
\end{verbatim}

\section{Extended Proofs}
\label{app:proofs}

\subsection{Proof of Proposition~\ref{prop:continuity} (Mixed Loss Convexity)}

\begin{proof}
We verify each axiom of a convex risk measure~\citep{follmer2002convex} for $\mathcal{L}_{\mathrm{mix}}(\alpha;\cdot) = \alpha\,\cvar_{0.95}(\cdot) + (1-\alpha)\,\rho_\lambda(\cdot)$.

\textit{Monotonicity.} If $X \le Y$ a.s., then $\cvar_{0.95}(-X) \ge \cvar_{0.95}(-Y)$ and $\rho_\lambda(-X) \ge \rho_\lambda(-Y)$, so $\mathcal{L}_{\mathrm{mix}}(\alpha;X) \ge \mathcal{L}_{\mathrm{mix}}(\alpha;Y)$.

\textit{Translation invariance.} For constant $c$: $\cvar_{0.95}(-(X+c)) = \cvar_{0.95}(-X) - c$ and $\rho_\lambda(-(X+c)) = \rho_\lambda(-X) - c$, hence $\mathcal{L}_{\mathrm{mix}}(\alpha;X+c) = \mathcal{L}_{\mathrm{mix}}(\alpha;X) - c$.

\textit{Convexity.} For $\lambda_0 \in [0,1]$ and P\&L's $X,Y$:
\begin{align}
    \mathcal{L}_{\mathrm{mix}}(\alpha;\lambda_0 X + (1-\lambda_0)Y)
    &= \alpha\,\cvar_{0.95}(-\lambda_0 X - (1-\lambda_0)Y) + (1-\alpha)\,\rho_\lambda(-\lambda_0 X - (1-\lambda_0)Y) \notag\\
    &\le \alpha[\lambda_0\,\cvar_{0.95}(-X) + (1-\lambda_0)\,\cvar_{0.95}(-Y)] \notag\\
    &\quad + (1-\alpha)[\lambda_0\,\rho_\lambda(-X) + (1-\lambda_0)\,\rho_\lambda(-Y)] \notag\\
    &= \lambda_0\,\mathcal{L}_{\mathrm{mix}}(\alpha;X) + (1-\lambda_0)\,\mathcal{L}_{\mathrm{mix}}(\alpha;Y),
\end{align}
where the inequality uses convexity of both $\cvar_{0.95}$ and $\rho_\lambda$.

Continuity in $\alpha$ follows from the linearity of the convex combination for fixed P\&L.
\end{proof}

\subsection{Proof of Proposition~\ref{prop:regret} (Regret Reduction)}

\begin{proof}
Let $\theta_t$ be the parameter at annealing step $t$ of Stage~2 with mixed gradient $g_t = \alpha_t g_{\mathrm{C}}(\theta_t) + (1-\alpha_t)g_{\mathrm{E}}(\theta_t)$.  Under $L$-smoothness of $\mathcal{L}_{\mathrm{mix}}$:
\begin{equation}
    \mathcal{L}_{\mathrm{mix}}(\alpha_{t+1};\theta_{t+1}) \le \mathcal{L}_{\mathrm{mix}}(\alpha_t;\theta_t) - \eta\|g_t\|^2 + \frac{L\eta^2}{2}\|g_t\|^2 + |\alpha_{t+1}-\alpha_t|\cdot C_\alpha,
\end{equation}
where $C_\alpha$ bounds $|\partial\mathcal{L}_{\mathrm{mix}}/\partial\alpha|$.  Summing over $K$ steps:
\begin{equation}
    \sum_{t=1}^K \eta(1 - L\eta/2)\|g_t\|^2 \le \mathcal{L}_{\mathrm{mix}}(\alpha_1;\theta_1) - \mathcal{L}_{\mathrm{mix}}(\alpha_K;\theta_K) + K\Delta\alpha\cdot C_\alpha.
\end{equation}
At the Stage~1 optimum $\theta_1^*$, we have $g_{\mathrm{C}}(\theta_1^*) \approx 0$, so $g_t \approx (1-\alpha_t)g_{\mathrm{E}}(\theta_1^*)$ initially.  The cumulative progress towards the entropic minimum is:
\begin{equation}
    \sum_{t=1}^K \eta(1-\alpha_t)^2\|g_{\mathrm{E}}(\theta_1^*)\|^2 \ge \frac{\eta K}{2}\|g_{\mathrm{E}}(\theta_1^*)\|^2 \cdot \overline{(1-\alpha)^2},
\end{equation}
where $\overline{(1-\alpha)^2} = \frac{1}{K}\sum_t(1-\alpha_t)^2$.  For $\alpha: 0.8\to 0.2$, $\overline{(1-\alpha)^2} \approx 0.37$.  This progress reduces the transition regret by $\frac{\eta K}{2}\|g_{\mathrm{E}}\|^2 \cdot 0.37$ relative to the two-stage protocol (which has $K=0$).
\end{proof}

\section{Numerical Stability}
\label{app:numerical}

\textbf{CVaR sorting.}  The empirical CVaR$_{0.95}$ is computed by sorting the batch P\&L and averaging the worst $\lfloor (1-0.95)B \rfloor$ samples.  For $B=256$, this uses the worst 12--13 samples, which can exhibit high gradient variance.  We mitigate this via the Rockafellar--Uryasev smooth surrogate~\eqref{def:cvar}.

\textbf{Mixed loss numerical conditioning.}  When $\alpha$ is close to 0 or 1, one component dominates and the gradient of the other may be numerically negligible.  Our schedule $\alpha \in [0.2, 0.8]$ avoids these extremes by construction.

\section{LSTM Gradient Flow and Delta Bounding}
\label{app:lstm_stability}

The LSTM architecture~\eqref{eq:lstm_f}--\eqref{eq:lstm_delta} mitigates the vanishing/exploding gradient problem through its gated cell state $c_k$.  The gradient of the loss w.r.t.\ $c_{k-1}$ factors as:
\begin{equation}
    \frac{\partial c_k}{\partial c_{k-1}} = f_k + \frac{\partial f_k}{\partial c_{k-1}} \odot c_{k-1} + \frac{\partial(i_k \odot \tilde{c}_k)}{\partial c_{k-1}}.
    \label{eq:cell_grad}
\end{equation}
When the forget gate $f_k \approx 1$ (retaining memory) and the remaining terms are small, $\partial c_k/\partial c_{k-1} \approx I$, allowing gradients to flow unattenuated across time steps.  This is critical for the 30-step hedging horizon ($N=30$), where vanilla RNNs suffer gradient decay of order $\sim 0.95^{30} \approx 0.21$.

The $\tanh$ delta bounding $\delta_k = \delta_{\max}\tanh(W_\delta h_k + b_\delta)$ ensures $|\delta_k| \le 1.5$ and provides a smooth gradient:
\begin{equation}
    \frac{\partial \delta_k}{\partial h_k} = \delta_{\max}(1 - \tanh^2(W_\delta h_k + b_\delta)) \cdot W_\delta,
\end{equation}
which is bounded by $\delta_{\max}\|W_\delta\|$ and vanishes only when $|W_\delta h_k + b_\delta| \to \infty$ (saturated regime).  The gradient clipping at $\|\nabla\| \le 5.0$ provides an additional safeguard.

\end{document}
